\scanpage{21}

les sous-a\deleted{\unsure{g|gg}}lgèbres égales à leur propre normalisateur).

Démonstration. a) Si \hspace{1.5em}$\operatorname{Nil}(a\deleted{\unsure{,|.}})\quad)$, comme cette dernière est nil\unsure{p}potente,
\hspace{1em}l'est. Inversement, si \hspace{1.1em} est nilpotent, elle est contenue dans
le nil-espace de son élément $a$, donc \hspace{2em}$\operatorname{Nil}(a\deleted{,})\quad)$.

b) Si \hspace{2em}$\operatorname{Nil}(a\deleted{,})\quad)$, comme $\operatorname{ad}(a)_{\quad/\operatorname{Nil}(a\deleted{,})}$ est injectif, il en
est à fortiori de $\operatorname{ad}(a)_{\quad/\quad}$, donc \hspace{1.5em} est son propre normalisateur
connexe. \deleted{\unsure{En|In}}Inversement, si \hspace{1em} est son propre normalisateur connexe, alors

% archive sheet number: 20
